\section{What a Finite Window Can Reproduce}
\label{sec:criterion}

A window can support accurate emulation only when histories with identical
recent inputs also produce nearly identical reservoir outputs. We quantify
the uncertainty left by the missing past, then establish when a sufficiently
expressive QELM can approach the best accuracy available from its window.

Consider histories with the same last $d$ inputs but potentially different
earlier inputs. Their largest output difference is
\begin{equation}
 D_d=\sup_{S_d(h)=S_d(h')}\|y(h)-y(h')\|.
 \label{eq:diameter}
\end{equation}
The supremum runs over all allowed histories. The norm is fixed and reduces
to absolute value for a scalar output. Thus $D_d$ measures how much the
output can vary when everything the window model receives is held fixed.

For a scalar output, each window is compatible with a range of reservoir
outputs. Predicting its midpoint minimizes the largest possible error:
any other prediction moves closer to one end but farther from the other.
If the range is $0.2$ to $0.8$, for example, predicting $0.5$ gives the
smallest worst-case error, $0.3$. Across all windows,
\begin{equation}
 \inf_f\sup_h |y(h)-f(S_d(h))|=\frac{D_d}{2},
 \label{eq:window-law}
\end{equation}
where $f$ ranges over all functions of the window. This limit applies
regardless of architecture or capacity. A longer window can reduce it by
distinguishing histories that a shorter window treats as identical. If
$D_d\to0$, arbitrarily small errors become possible in principle. If the
gap does not vanish, additional capacity cannot remove the error floor.

Approaching this limit with a QELM requires suitable dynamics and model
capacity. Assume a finite-dimensional reservoir, a compact input alphabet,
input-continuous quantum channels, and fixed observable outputs. Every
finite input continuation must be allowed after every admissible earlier
history. These conditions make the midpoint prediction continuous in
the window.

For each window length, let $\mathcal Q_d$ be a reset-QELM family that can
approximate every continuous window function uniformly as capacity grows,
using exact expectation values. Then, for scalar outputs,
\begin{equation}
 \inf_{f\in\mathcal Q_d}\sup_h
 |y(h)-f(S_d(h))|=\frac{D_d}{2}.
 \label{eq:qrc-qelm-equivalence}
\end{equation}
The family can approach the window's information limit arbitrarily closely.
This is an infimum: a particular finite model need not attain it exactly.

Consequently, under these assumptions, arbitrarily accurate QELM emulation
is possible precisely when $D_d\to0$ as the window grows. The omitted
past must become negligible uniformly across all allowed histories, not
only typical test examples. The same criterion holds for finite vector
outputs, although their unrestricted fixed-window optimum lies between
$D_d/2$ and $D_d$. Appendix~\ref{app:window-laws} provides the proofs and
shows that randomization cannot lower the scalar bound on worst-case
expected absolute error.

For inputs in $[0,1]$, repeated encoding and Bernstein approximation give
an explicit growing QELM construction
\cite{goto2021universal,yu2024nonasymptotic}. Its resource costs are discussed
next. This existence result neither guarantees efficient training nor
makes the fixed architectures in Section~\ref{sec:numerics} universal.
The replacement question is now one of memory: when do older inputs lose
their influence, and how much context makes their omission acceptable?
